% 水平圆盘两侧物体（重心不在圆心）：T、f_1、f_2 随 \omega^2
% 横轴四刻度均匀；刻度旁为 \omega^2 公式。纵轴为力（取 \mu g=1,\ m_1=2,\ m_2=1,\ r_1=2,\ r_2=1 计算曲线），刻度旁为符号式。
% x\in[0,1]\to u\in[0,1],\ [1,2]\to[1,4/3],\ [2,3]\to[4/3,2]，其中 u=\omega^2 r_1/(\mu g)。
\begin{center}
  \footnotesize
  % \begin{tikzpicture}
  %   \pgfmathsetmacro{\xmin}{-0.12}
  %   \pgfmathsetmacro{\xmax}{3.12}
  %   \begin{axis}[
  %     width=0.94\columnwidth,
  %     height=6.4cm,
  %     xmin=\xmin,
  %     xmax=\xmax,
  %     ymin=-1.18,
  %     ymax=2.38,
  %     axis lines=middle,
  %     x axis line style={-stealth},
  %     y axis line style={-stealth},
  %     xlabel={$\omega^2$},
  %     xlabel near ticks,
  %     ylabel={$T,\ f_1,\ f_2$},
  %     ylabel near ticks,
  %     xtick={0,1,2,3},
  %     xticklabels={%
  %       $0$,
  %       $\dfrac{\mu g}{r_1}$,
  %       $\dfrac{\mu m_1 g}{m_1 r_1-m_2 r_2}$,
  %       $\dfrac{\mu (m_1+m_2) g}{m_1 r_1-m_2 r_2}$%
  %     },
  %     xticklabel style={align=center, font=\scriptsize},
  %     ytick={-1,0,0.5,2/3,2},
  %     yticklabels={%
  %       $-\mu m_2 g$,
  %       $0$,
  %       $\dfrac{\mu m_2 g r_2}{r_1}$,
  %       $\dfrac{\mu m_1 g\, m_2 r_2}{m_1 r_1-m_2 r_2}$,
  %       $\mu m_1 g$%
  %     },
  %     yticklabel style={font=\scriptsize, align=right},
  %     legend style={
  %       at={(0.5,1.02)},
  %       anchor=south,
  %       legend columns=-1,
  %       /tikz/every even column/.append style={column sep=0.5em},
  %     },
  %     legend cell align=left,
  %     clip=false,
  %   ]
  %     % 竖直虚线：三临界
  %     \addplot [dashed, gray!70, forget plot] coordinates {(1,-1.18) (1,2.38)};
  %     \addplot [dashed, gray!70, forget plot] coordinates {(2,-1.18) (2,2.38)};
  %     \addplot [dashed, gray!70, forget plot] coordinates {(3,-1.18) (3,2.38)};

  %     % 水平虚线：特殊力值
  %     \addplot [dashed, gray!55, forget plot] coordinates {(\xmin,-1) (\xmax,-1)};
  %     \addplot [dashed, gray!55, forget plot] coordinates {(\xmin,0.5) (\xmax,0.5)};
  %     \addplot [dashed, gray!55, forget plot] coordinates {(\xmin,2/3) (\xmax,2/3)};
  %     \addplot [dashed, gray!55, forget plot] coordinates {(\xmin,2) (\xmax,2)};

  %     % T（\mu g=1）：x\in[0,1] 为 0；x\in[1,3] 上 u 分段线性，T=2(u-1)
  %     \addplot [very thick, black] coordinates {
  %       (0,0) (1,0) (2,2/3) (3,2)
  %     };
  %     \addlegendentry{$T$}

  %     % f_1
  %     \addplot [very thick, red!75!black] coordinates {
  %       (0,0) (1,2) (3,2)
  %     };
  %     \addlegendentry{$f_1$}

  %     % f_2
  %     \addplot [very thick, blue!75!black] coordinates {
  %       (0,0) (1,0.5) (2,0) (3,-1)
  %     };
  %     \addlegendentry{$f_2$}
  %   \end{axis}
  % \end{tikzpicture}\\[0.4em]
  \begingroup
  \rowcolors{1}{white}{white}%
  \begin{tabular}{|p{0.35\columnwidth}|p{0.1\columnwidth}|p{0.3\columnwidth}|}
    \hline
    阶段 & 拉力 $T$ & 摩擦力 $f_1,\ f_2$ \\
    \hline
    $\omega^2\leq \dfrac{\mu g}{r_1}$ &
      无 &
      两物体摩擦力均指向圆心，随 $\omega^2$ 增大 \\
    \hline
    $\dfrac{\mu g}{r_1}<\omega^2\leq \dfrac{\mu m_1 g}{m_1 r_1-m_2 r_2}$ &
      出现并增大 &
      $f_1$ 已达最大并保持；$f_2$ 先增至 $\dfrac{\mu m_2 g\,r_2}{r_1}$ 再减至 $0$ \\
    \hline
    $\dfrac{\mu m_1 g}{m_1 r_1-m_2 r_2}<\omega^2\leq \dfrac{\mu(m_1+m_2)g}{m_1 r_1-m_2 r_2}$ &
      继续增大 &
      $f_1$ 保持最大；$f_2$ 背离圆心增大至 $\mu m_2 g$ \\
    \hline
    $\omega^2>\dfrac{\mu(m_1+m_2)g}{m_1 r_1-m_2 r_2}$ &
      \multicolumn{2}{p{0.58\columnwidth}|}{无法保持相对静止，两物体相对圆盘滑动} \\
    \hline
  \end{tabular}
  \endgroup
\end{center}
